\section{Extended implications for human play}
\label{app:human}

\vfast{} maintains a $54 \times 6$ posterior and re-scores its candidate asks
after every public event. A person at a table cannot. The mechanisms that the
frozen-policy ablations of \S\ref{app:ablations} identify as the ones this
configuration depends on most heavily are, however, not the expensive ones, and
several of them can be run by hand. This appendix records them at length; the
condensed version is in \S\ref{sec:humanplay}.

\subsection{How to read the labels}
\label{app:human-labels}

Each item below carries one of three tags, because the recommendations differ
sharply in how well they are supported.

\begin{description}[leftmargin=2.2em,style=nextline]
\item[\textbf{[R]}] A logical consequence of the rules dialect of
\S\ref{sec:rules}. It holds for any opponent and requires no experiment. A tag
of \textbf{[R]} is never attached to a prescription about how to play; where a
rules consequence suggests a course of action, the action is split off and
tagged separately.
\item[\textbf{[M]}] A measured result. The experiment is named. In
every case the measurement is of \emph{this} fitted configuration on a specific
bank against a specific panel, and in the case of an ablation it is a
frozen-policy effect rather than the standalone value of the component.
\item[\textbf{[U]}] An untested suggestion. It is consistent with
the design of the agent or with the rules, but no experiment in this paper
isolates it, and it should be treated as a hypothesis.
\end{description}

\subsection{What an ask tells the table}
\label{app:human-ask}

\textbf{[R] An ask is a confession about the asker.} Under this dialect a
player may ask for a card only in a half-suit in which they already hold at
least one card. So a player who asks for the queen of hearts has stated two
things publicly: that they do not hold the queen of hearts, and that they hold
at least one other card of the high hearts. The second is the one that is easy
to miss. It applies to a teammate's ask exactly as much as to an opponent's.

\textbf{[R] A miss is a second confession, about the target.} A failed ask
establishes that the named card was not in the target's hand at that moment. The
exclusion stands until a public transfer moves that card, and every transfer in
this game is public, so the exclusion is never silently invalidated.

\textbf{[R] The content of both events is an exclusion, not an impression.}
What an ask and a miss add to the public record is a collection of statements of the
form ``this card is not held by this player'' together with a disjunction over
the asker's own hand. ``Somebody wanted low clubs'' is not among them; ``the
seven of clubs is not held by seat 2 or seat 4, and seat 3 holds at least one of
the other five low clubs'' is.

\textbf{[U] Write them down in that form.} Exclusions and disjunctions combine
with hand counts, and impressions do not, so a person keeping a record by hand
is better served by the first. No experiment in this paper evaluates any human
bookkeeping scheme.

\textbf{[M] The certificates that encode these two deductions are removed by the
largest frozen-policy ablations in the study.} In E5, dropping the ask-legality
certificates --- while keeping capacity conditioning, keeping every exclusion
from a miss, and keeping the entire decision policy unchanged --- moved the
frozen policy from \numAblRef\% to \numAblSinkhornAbs\%. Dropping capacity
conditioning as well moved it further, to \numAblCFourAbs\%. These are the
largest decreases in the study, but none of them isolates a single constraint:
the implementation applies the fitted policy prior only along the path these
rows replace, so each removes the prior too (\S\ref{sec:ablations}).
These are effects on this
frozen parameter vector; they do not say what a policy refitted without those
constraints would score.

\subsection{Counting}
\label{app:human-count}

\textbf{[R] Hand sizes are public, and they close deductions that exclusions
alone leave open.} If a player holds four cards and you have already located four
specific cards in their hand, every other unresolved card is excluded from them
at once. Conversely, if a player holds four cards and only four unresolved cards
can possibly be theirs, they hold all four.

\textbf{[U] Run the check after every transfer, on the two players whose
counts changed.} That is where a new certainty is most likely to appear. The
recommendation about \emph{when} to run the check is a suggestion about human
bookkeeping; no experiment here evaluates it.

\subsection{Locked half-suits, and what waiting does and does not buy}
\label{app:human-locked}

\textbf{[R] If your team holds all six cards of a half-suit, no opponent can
take any of them.} An opponent needs a card of a half-suit to ask in it, and an
opponent has none. Ownership of that half-suit is therefore frozen
(Theorem~\ref{thm:locked}), and an opponent who claims it must claim it
incorrectly, which awards it to you.

\textbf{[R] Waiting therefore carries no risk to ownership}
(Corollary~\ref{cor:wait}). No sequence of opponent asks can remove a card of a
half-suit your side holds in full, so the passage of time cannot cost your team
that half-suit.

\textbf{[U] Spend the intervening turns resolving the split.} If you are
certain your side holds all six but unsure which teammate holds what, there is
no time pressure from the opponents. Whether turns are better spent that way
than on anything else available is not measured here.

\textbf{[R] The rules push declaration timing in both directions.} One
consequence favours holding. A locked half-suit scores only when it is claimed,
and claiming resolves six cards that were, from an opponent's point of view,
part of the unresolved pool $U$; by \eqref{eq:capacity} that tightens every
opponent's capacity constraint on the cards that remain, so holding the
half-suit keeps six slots absorbing probability mass the opponents cannot
allocate elsewhere. The other favours claiming. A card of a locked half-suit
licenses only asks that are certain to fail, so a hand composed largely of
locked cards yields little from a turn, whereas a player who claims and empties
a hand leaves the asking cycle and passes the turn to a teammate whose asks can
succeed. A third channel runs alongside both: asks inside a
locked half-suit remain legal for the owning team and still carry ask-legality
certificates, so information about the split continues to change while ownership
does not (\S\ref{sec:termination}). Develin \cite{develin} records that
channel from the human side --- such asks are available to the owning team and
to nobody else --- and the fitted \vfast{} configuration was given no term that
uses it (\S\ref{sec:limitations}).

\textbf{[U] Which of the two dominates is a question about the position in
front of you.} It is quantitative, it depends on the opponents, and
Theorem~\ref{thm:locked} does not settle it. Nothing in this paper measures the
trade-off directly.

\textbf{[R] If a half-suit is not yet fully yours, the ownership guarantee does
not apply.} An opponent holding a card of the half-suit has a legal ask in it,
so a card you hold can still be taken away.

\textbf{[U] In that case prefer the earlier claim.} Weigh the risk that an
opponent removes one of your cards against the chance that the split is not yet
resolved. No experiment here isolates the value of claiming earlier in an
unlocked half-suit.

\textbf{[R] A hand made entirely of locked cards licenses only asks that are
certain to fail.} Every possible target of such an ask is void in the half-suit,
so the ask transfers nothing and the turn produces no card movement.

\textbf{[U] So claim the half-suit, empty the hand, and pass the turn to a
teammate who can use it.} That is a suggestion about what to do with a turn
that cannot produce a transfer; the rules consequence above is what is
established, and the fitted configuration was not given a term that targets this
case.

\textbf{[M] No experiment here isolates the value of waiting.} The
``cash locked half-suits immediately'' switch is inert in the shipped
configuration --- the value-based stopping rule decides first, so the switch is
never consulted --- and its E5 row is exactly zero by construction. What is
observed, not manipulated, is that on the primary bank the fitted \vfast{}
configuration held a locked half-suit for \numLockHoldFour{} public events
before a correct declaration, against \numLockHoldThree{} for \vthree{}
(\S\ref{sec:results}); that is a description of two policies' behaviour, not
evidence that either duration is better. What the rules guarantee is the safety
of waiting with respect to ownership, not its profit.

\subsection{Ranking candidate asks}
\label{app:human-rank}

The ordering below is roughly what the fitted weights produce. The support
behind the individual steps is very uneven, and the tags say so. Every interval
quoted in this subsection is a paired percentile bootstrap over deals, the
construction described in \S\ref{app:abl-design}.

\begin{enumerate}[leftmargin=2em]
\item \textbf{[R] An ask for a card you have located cannot fail.} The card is
in the named hand, so the ask succeeds and the turn stays with you.
\textbf{[U] Take every such ask before anything speculative.} The ordering is a
suggestion; no experiment here compares it against alternatives.
\item \textbf{[R] A hit retains the turn, so an ask can start a sequence.} The
value of an ask is therefore not confined to the single card it names.
\textbf{[U] Prefer a likely card that leaves another likely card behind it over
an equally likely card that leaves nothing.} No experiment here isolates this
ordering.
\item \textbf{[M] Then the ask that completes a lock.} An ask that would give
your team the sixth card of a half-suit converts a contested half-suit into a
frozen one. In E5, zeroing the lock-completion weight cost \numAblLockWeight{}
points (95\% paired percentile-bootstrap interval \numAblLockWeightCI{}) on the
frozen configuration --- one of the few individual ask weights whose interval
excludes zero.
\item \textbf{[U] Then consider who receives the turn on a miss.} A failed ask
hands the turn to the player you named, so you are choosing a successor as much
as a card; avoid handing it to a player visibly strong in a half-suit where your
side is exposed. In E5, zeroing the reply-threat weight changed the win rate by
\numAblThreat{} points with a paired interval of \numAblThreatCI{}, so this
study does not establish the effect.
\item \textbf{[U] Then consider what the ask gives away.} Your first ask in a
half-suit announces that you hold one of its cards. Between two near-equal asks,
the one in a half-suit you have already revealed leaks less. In E5, zeroing both
information-leak weights changed the win rate by \numAblLeak{} points with a
paired interval of \numAblLeakCI{}; again not established.
\end{enumerate}

\subsection{Checking a declaration before making it}
\label{app:human-declare}

\textbf{[R] Two different claims are involved, and only the stronger is
sufficient.} ``Our side holds all six of these cards'' and ``I know which of us
holds each one'' are not the same statement: a declaration names a
card-to-player allocation, and an allocation that is right about the team and
wrong about the split fails.

\textbf{[U] Before claiming, name all six cards and a holder for each,
explicitly.} If the team-level claim is solid but the split is not, and the
half-suit is already locked, wait and resolve the split. If it is not locked,
weigh the incomplete claim against the chance that an opponent removes one of
your cards first. The procedure is a suggestion about how a person might make the
check; nothing in this paper evaluates it.

\textbf{[M] The measured difference between the two agents is in declaration
accuracy; which half of the mechanism produces it is not resolved.} On the
primary bank \vfast{} declared correctly \numVsVthreeDeclAcc\% of the time
against \vthree{}'s \numVthreeDeclAccSame\% in the same games
(\S\ref{sec:results}). In E5, replacing the fitted stopping rule with a fixed
confidence threshold changed the win rate by \numAblStop{} points (95\% paired
percentile-bootstrap interval \numAblStopCI{}), and moving the declaration
confidence threshold between $0.80$ and $0.99$ also changes nothing the
experiment can resolve. That stopping-rule ablation is
unresolved from zero, so these experiments do not attribute the
declaration-accuracy difference to the probability estimate rather than to the
rule that consumes it. The calibration detail in \S\ref{app:calibration} is a
description of that estimate, not an explanation of the behaviour.

\subsection{The forced endgame}
\label{app:human-endgame}

\textbf{[R] Running out of cards is not a defeat.} A player with no cards cannot
be asked, so nothing can be taken from them. A whole team with no cards forces
the opponents to claim every remaining half-suit without asking, and some of
those claims will be wrong, which awards those half-suits to the cardless team.

\textbf{[U] Emptying a hand deliberately is an option worth knowing, not a
measured recommendation.} The preceding paragraph follows from the rules; no
experiment in this paper isolates the value of steering towards a cardless hand
on purpose, and the fitted configuration was not given a term that rewards it.

\textbf{[R] A correct claim is informative to everyone; an incorrect one is
informative and costly.} A correct claim announces the exact split of six cards,
which changes every player's count of what remains and can resolve a half-suit
that was previously unresolved. An incorrect claim supplies the same information
to the opponents and gives them the half-suit as well.

\textbf{[U] When it is your team that must claim everything, claim the surest
half-suit first.} Ordering the claims by confidence makes the later claims
easier, on the argument above; the reverse order forgoes that help. The
engine's forced endgame uses a fixed willingness ladder rather than a confidence
ranking (\S\ref{sec:declare}), and no experiment here compares claim orders.

\subsection{Deliberate misdirection}
\label{app:human-bluff}

\textbf{[M] Nothing in this study supports paying for theatre.} The
v0.3 study reported no reliable benefit from deliberate misdirection, and this
study does not change that. \vfast{} contains no bluffing rule, and on the
primary bank it won \numVsBluffer\% of games against the misdirection artist,
the archetype in the panel built around deceptive asking
(Table~\ref{tab:h2h}). That is the second largest margin in the panel, behind
the random legal control at \numVsRandom\% and ahead of the focused hunter at
\numVsHunter\%, the adaptive diversifier at \numVsDiversifier\% and \vtwo{} at
\numVsVtwo\%: the deception archetype is not separated from the panel's other
weak baselines.

That is a result about one archetype in one panel, and it is not an isolation
experiment: the misdirection artist differs from the other baselines in more
than its use of deception, so the margin cannot be attributed to the deception
alone. The rules-level observation behind the negative result is narrower and
firmer: an ask is legal only from a hand that holds another card of the
half-suit, so the information an ask reveals is constrained by what the asker
actually holds. Behaviour that changes neither what the other players can deduce
nor who holds the turn changes nothing that the game scores.
